\scanpage{309}
Questions sur la structure de $\uline{\mathrm{Pic}}_{X/S}$ \quad ($S$ noeth.)

\margin{O.K. $\Longleftarrow$}
a) $X$ est-il somme \inserted{(disjointe)} de préschémas de type fini sur $S$ ?
\margin{non en général; oui ``géométriquement sur $S$''.}

\margin{$\Longleftarrow$}
\unsure{a'}\deleted{\unsure{b}}) Pour tte partie quasi-compacte $U$ de $X$, est-il vrai que $\overline U$ est q. cpt ? $\Longleftrightarrow$ Pour tt $x\in X$, est-il vrai que $\overline x$ est quasi-compact ?

a'') Est-il vrai que \deleted{\unsure{le schéma}} $\uline{\mathrm{Pic}}^\tau_{X/S}$ est un sous-préschéma \uline{fermé} ?? (Il est en \uline{tous cas \unsure{ouvert}})
\margin{?}

On \unsure{aurait} $a\Rightarrow a'\Rightarrow a''$. La réponse à $a''$ est affirmative si $X$ est \deleted{\unsure{simple}}\inserted{\unsure{normal}} (et \unsure{propre}) sur $S$, mais je ne sais pas pour autant s'il en est de $\unsure{\overline{=}}$ pour $a'$).
\margin{\unsure{l'équiv} au th. de décomposition (résulte du th. d'existence en Géom. formelle); il suffit en \unsure{fait} \unsure{que} un Pic. \unsure{connexe} \unsure{des fibres en soit complet}.}

b) Question: existence du préschéma quotient
\[
\uline{\mathrm{N\acute{e}ron}}^0_{X/S}\inserted{\underset{}{\unsure{\mathrm{NSL}}}}=\uline{\mathrm{Pic}}_{X/S}/\uline{\mathrm{Pic}}^\tau_{X/S}
\]

c) $\uline{\mathrm{Pic}}^\tau_{X/S}$ est-il universellement ouvert sur $S$ ?
\margin{O.K. en les pts \unsure{correspondant} à une torsion de N.S. première à la car. résiduelle}
\margin{\uline{Plat sur $S$} ?}

d) $\uline{\mathrm{Pic}}^\tau_{X/S}$ est-il quasi-compact \inserted{donc de type fini} sur $S$ ?
\margin{\uline{Oui, toujours}}
\margin{(O.K. si $X$ normal sur $S$, \unsure{via} \unsure{un} ``critère d'\unsure{irrégularité|irréductibilité}'' de Weil, à démontrer via \unsure{Géom. formelle} (via lemme Zariski); ? ? ?}
\margin{(OK. si les fibres sont \unsure{primaires|prénormales})}

e) $\uline{\mathrm{Pic}}^\tau_{X/S}$ est-il séparé sur $S$ ?
\margin{\uline{non} si on ne le suppose pas \unsure{représentable}; l'est-il s'il est représentable ?}

\deleted{\unsure{f}}\inserted{\unsure{b'}}) demi-continuité de (\unsure{Min.|Dim.} Lissité) des fibres ....

\margin{$\leftarrow$}
\unsure{f}) Continuité de $\unsure{\ell}$-\unsure{rang} des Picard des fibres
\margin{OK si $X$ normal sur $S$, plus gén. si les \unsure{comp. connexes} de Picard des fibres sont complètes, via groupe fondamental}

g) $\uline{\mathrm{Pic}}^\tau\deleted{\unsure{(X\times_S Y)}}_{X\times Y/S}\simeq\uline{\mathrm{Pic}}^\tau_{X/S}\times_S\uline{\mathrm{Pic}}^\tau_{Y/S}$
\margin{O.K. \deleted{si $X$ et $Y$ sont à fibres \unsure{géom. intègres|géom. connexes} [condition plus ou moins \unsure{nécessaire} vérifier.]}}

h) Si $Y$ section hyperplane de $X$,
\[
\uline{\mathrm{Pic}}^\tau_{X/\unsure{k|u}}\longrightarrow\uline{\mathrm{Pic}}^\tau_{Y/\unsure{k|u}}\quad\text{bijectif ? \unsure{Isom?|Sous?}}
\]
\[
\uline{\mathrm{Pic}}_{X/\unsure{k|u}}\mathbin{\unsure{-|\to}}\uline{\mathrm{Pic}}^{\unsure{\tau|0}}_{Y/\unsure{k|u}}\quad\text{\unsure{surjectif}}
\]
\margin{conditions à préciser....}

\figure{\begin{tikzpicture}[x=1cm,y=1cm]
\draw (0,4) .. controls (-.4,3) and (-.65,1) .. (0,0);
\draw (-.1,4.05) .. controls (0,3.85) and (.12,3.85) .. (.2,4.02);
\draw (0,0) -- (-.05,.3);
\end{tikzpicture}}
\figure{\begin{tikzpicture}[x=1cm,y=1cm]
\draw (0,5.8) .. controls (-.2,6.2) and (1.1,6.4) .. (1.4,6.05) .. controls (1.6,5.75) and (.3,5.55) .. (0,5.8);
\draw (.15,5.72) -- (.1,5.4);
\draw (2.9,4.25) .. controls (2.4,4.45) and (2.6,5.15) .. (3.15,5.12) .. controls (3.8,5.3) and (4.6,4.9) .. (4.85,4.86);
\draw (2.85,4.3) .. controls (3.4,4.05) and (4.3,3.95) .. (4.9,4.05);
\draw (2.82,4.7) .. controls (2.3,4.35) and (2.2,4.1) .. (2.05,3.8);
\draw (0,2.65) .. controls (-.3,3.05) and (1.9,3.1) .. (2.1,2.8) .. controls (2.2,2.45) and (.15,2.4) .. (0,2.65);
\draw (.15,2.5) -- (0,2.28);
\draw (2.75,.1) -- (3.65,2.15);
\draw (3.15,0) -- (4.05,2.05);
\draw (3.7,0) -- (4.5,1.95);
\draw (4.25,.08) -- (4.85,1.6);
\end{tikzpicture}}
\mspage{308}
